%% 01-introduction.tex -- generated by tools/build_from_docx.py; edit freely

\hypertarget{introduction}{%
\section{Introduction}\label{introduction}}

Sixth-generation networks are expected to be managed at the timescales
on which they change, by intelligence placed throughout the
infrastructure rather than concentrated above it. In UNITY-6G that
intelligence is not supplied by one party. The algorithms and
applications that observe, analyze, decide and act are contributed by
different partners, address different domains --- the core, the open
radio access network, non-terrestrial segments, wireless transport,
non-3GPP access, the AI/ML domain, vertical applications and digital
twins --- and are realized with techniques that range from rules and
solvers through trained models to language-model reasoning. They
nevertheless act on shared infrastructure, often on the same resources,
often toward objectives that pull against each other, and often at the
same moment.

The unit in which this intelligence is hosted is the agent: a
composition of four capabilities, instantiated against a declared
objective, described by a signed descriptor, and holding only the
authority the platform grants it. The agentic framework is the set of
contracts and platform obligations under which agents are admitted,
placed, bound to one another and to the infrastructure, mediated,
executed, verified and retired. Its contracts are stated precisely
enough to be implemented without further interpretation: every
declaration has a schema, every check a point at which it runs, and
every obligation a party that carries it.

\hypertarget{authority-as-a-shared-resource}{%
\subsection{Authority as a Shared Resource}\label{authority-as-a-shared-resource}}

Making a single management loop intelligent is not the difficult part.
The difficult part is hosting many of them, independently authored, over
one infrastructure, and four questions arise that a single loop never
has to answer. Who may act on what, and how is that settled before
anything acts? What happens when two agents, each reasoning correctly
within its own domain, decide on actions that conflict, when the
information needed to see the conflict exists in neither of them? How
does anyone know that an action took effect, rather than that a
controller acknowledged it? And how does the party that requested a
service learn whether the service is being delivered, rather than
whether its agents are busy?

The framework answers each question with a mechanism that sits outside
the agents and does not depend on their good behavior. Authority is
allocated by the platform at admission, bounded by a declared scope and
action space, held in a registry and revocable at any time; an agent
never grants itself authority and never approves its own decision.
Intents that could conflict are mediated by a conflict fabric before any
of them reaches the network. Actuation is verified against independently
observed state, and an outcome that cannot be verified is reported as
unverified. Assurance is assembled from events the platform emits rather
than from what agents report about themselves. Authority to change the
network is, in short, treated as a shared resource: granted, bounded,
attributed and withdrawn by the platform.

\hypertarget{the-agent}{%
\subsection{The Agent}\label{the-agent}}

Each of the four capabilities of an agent is defined by what it produces
and, as importantly, by what it never produces. The Monitoring System
produces evidence and never conclusions. The Analytics Engine produces
insight and never action. The Decision Engine produces intent, never a
controller command, and never approves its own intent. The Actuator
executes mediated intents through a controller, verifies the result and
reports one of four typed outcomes; actuation is a role, not a
privilege. A capability may be realized by the agent in full or in part,
delegated to a provider, or absent, and the composition that results
sets the highest authority the agent can hold.

An agent is local when it acts within one domain under that domain's
orchestrator, and global when it operates at the inter-domain level,
where it acts only by delegating objectives and bounds to local agents.
What an agent is, what it may touch, how fast it runs, where it may be
placed, how it fails and who answers for it are declared in its
descriptor, which is the single contract between the contributor and the
platform. The framework is flexible in what may be declared and strict
in requiring that it be declared, checked and adhered to, and the same
contracts hold whatever technique realizes a capability, because the
bounding comes from the platform rather than from the technique.
Responsibility for an agent's behavior rests with the partner who
contributed it, and the platform keeps the evidence that makes every
transaction attributable to that partner and to the descriptor under
which the agent was admitted.

\hypertarget{place-in-the-unity-6g-architecture}{%
\subsection{Place in the UNITY-6G Architecture}\label{place-in-the-unity-6g-architecture}}

Agents live inside the hierarchical management and orchestration of
UNITY-6G. The Inter-Domain Management Orchestrator (IDMO) decomposes a
customer-facing service into resource-facing service requests, one per
domain, and that pipeline realizes services rather than agents. An agent
is instantiated from its descriptor, which carries the service tree of
its tuples: a global agent is issued to the IDMO, which resolves the
tree into one portion per domain, and a local agent is issued directly
to its Domain Management Orchestrator (DMO). The DMO realizes the
portion that falls to it, and every DMO wraps the platform native to its
domain --- the RAN intelligent controllers of the open radio access
network, the software-defined networking controller of the transport
domain, the network function orchestration of the core --- rather than
replacing it. Global agents are held at the IDMO; local agents are
placed, bound and supervised by their DMO.

Two communication planes serve the agents, and they are kept apart. The
service-based management bus (SMB) carries control and coordination
only: registration, lifecycle, intents, outcomes and compact insight.
Observational data travels the Data Continuum, as governed streams and
data products, and never the SMB. A policy that lets an agent act
without a round trip is requested by the agent from the party that
grants it, and is never pushed to an agent that did not request it.

The timescales on which agents operate are the latency classes of the
open radio access network --- real time below 10 ms, near real time from
10 ms to 1 s, and non-real time at 1 s and above --- together with an
offline class for agents that never act on a live loop. The management
vocabulary is that of 3GPP: intent-driven management, closed-loop
automation, service lifecycle states and service-level assurance. The
framework binds to that vocabulary and to the standard interfaces of
both bodies rather than defining alternatives to them.

Agent-based services are not the only services in the architecture.
Services managed conventionally, without agents, coexist with them in
the same deployment and are unaffected by the framework, which governs
what the agentic construct requires wherever it is used.

\hypertarget{division-of-responsibility}{%
\subsection{Division of Responsibility}\label{division-of-responsibility}}

The framework binds several parties, each of which owns part of what an
agent depends on. T2.3 owns the agent contracts, the descriptor, the
platform obligations toward agents, and the catalog entries specific to
the framework. Conflict resolution algorithms and trust modeling are
owned by WP4 and enter the platform as services; the framework states
only the surface through which agents meet them, so that WP4 decides and
the framework binds. Objective KPIs and their measurement methodology
are owned by T2.1. The metrics, control parameters, action types and
controller bindings of a domain are owned by that domain. Partners
contributing agents own the behavior of what they contribute. Table 1-1
states what each party owns and how the framework relates to it.

\begin{wptable}{1-1}

\begin{longtable}[]{@{}
  >{\raggedright\arraybackslash}p{(\columnwidth - 4\tabcolsep) * \real{0.2125}}
  >{\raggedright\arraybackslash}p{(\columnwidth - 4\tabcolsep) * \real{0.4750}}
  >{\raggedright\arraybackslash}p{(\columnwidth - 4\tabcolsep) * \real{0.3126}}@{}}
\caption{Parties bound by the framework and what each
owns}\tabularnewline
\toprule\noalign{}
\begin{minipage}[b]{\linewidth}\raggedright
\textbf{Party}
\end{minipage} & \begin{minipage}[b]{\linewidth}\raggedright
\textbf{Owns}
\end{minipage} & \begin{minipage}[b]{\linewidth}\raggedright
\textbf{Relation to the framework}
\end{minipage} \\
\midrule\noalign{}
\endfirsthead
\toprule\noalign{}
\begin{minipage}[b]{\linewidth}\raggedright
\textbf{Party}
\end{minipage} & \begin{minipage}[b]{\linewidth}\raggedright
\textbf{Owns}
\end{minipage} & \begin{minipage}[b]{\linewidth}\raggedright
\textbf{Relation to the framework}
\end{minipage} \\
\midrule\noalign{}
\endhead
\bottomrule\noalign{}
\endlastfoot
\textbf{T2.3} & The agent contracts and the descriptor schema; the
admission, activation, lifecycle and assurance obligations; the
agent-side surface of mediation; the reason codes and verification
methods of the type catalog. & Defines them. \\
\textbf{WP4} & Conflict resolution algorithms and the conflict fabric;
trust modeling. Both enter as platform services. & Bound through the
four-interaction surface of Chapter 10, and consulted at admission and
registration for trust. \\
\textbf{T2.1} & Objective KPIs, their units and their measurement
methodology. & Referenced through the objective KPI entries of the type
catalog. \\
\textbf{IDMO} & Decomposition of customer-facing services, delegations
to domains, and the global agents. & Carries the obligations of the
inter-domain level. \\
\textbf{DMOs and domain owners} & Realization and placement within the
domain; the domain\textquotesingle s metrics, control parameters, action
types, controller bindings and source kinds. & Carry the obligations of
the domain level and own the catalog entries of their domain. \\
\textbf{Contributing partners} & The logic of the agents they
contribute, and responsibility for its behavior. & Held to the role
contracts at admission and in operation; every transaction is
attributable to them. \\
\end{longtable}

\end{wptable}
